% 第16批：原书印刷第94—95页；源图题标注复核条件。
\subsection{自测题4：简谐振动、波形及驻波填空}

\sourcetest{4.2.1}{94}{102}
\problemtext 频率 $0.25$Hz，$t=0$ 时位移为 $-0.37$cm 且速度为零；求振幅及余弦型振动式。
\workings 零速即转折点，$A=0.37$cm，$\omega=2\pi f=\pi/2$rad/s。故 $\boxed{x=0.37\cos(\pi t/2+\pi)\ {\rm cm}}$。
\crosscheck 初始在负端点，必须取初相 $\pi$，不是0。

\sourcetest{4.2.2}{95}{103}
\problemtext 竖直弹簧振子位移为振幅的一半时，求动能与总能量的比，以及由静伸长 $\Delta l$ 表示的周期。
\workings $E=\frac12kA^2$；位移为 $A/2$ 时 $K/E=1-(1/2)^2=\boxed{3/4}$。静平衡 $k\Delta l=mg$，$T=2\pi\sqrt{m/k}=\boxed{2\pi\sqrt{\Delta l/g}}$。
\crosscheck 势能零点取弹簧振子的重力修正平衡位置。

\sourcetest{4.2.3}{95}{103}
\problemtext 题4.2.1图示振幅为4个纵轴单位，$t=0$ 时 $x=-2$ 且正在向正向运动，$t=2$s 下降经过平衡点；写出周期和余弦式初相。
\workings 下降过平衡点时的相位为 $\pi/2$（模 $2\pi$）；$t=0$ 为 $\cos\varphi=-1/2$ 且 $v>0$，可取 $\varphi=-2\pi/3$。从 $-2\pi/3$ 到 $\pi/2$ 的相位跨度为 $7\pi/6$，据此 $\omega=7\pi/12$rad/s、$\boxed{T=24/7\ {\rm s}}$。
\crosscheck 图中零点的时间标注需要独立对照：若印刷图仅为示意而非精确刻度，周期不应由图形长度估读；本条保留待复核。

\sourcetest{4.2.4}{95}{103}
\problemtext 同向同频简谐振动合振幅为20cm，第二振动振幅 $10\sqrt3$cm，且 $\varphi_2-\varphi_1=\pi/6$；求第一振幅及逆序相位差。
\workings $20^2=A_1^2+(10\sqrt3)^2+2A_1(10\sqrt3)\cos(\pi/6)$，得 $A_1^2+30A_1-100=0$。取正根 $\boxed{A_1=5\sqrt{13}-15\simeq3.028\ {\rm cm}}$；$\boxed{\varphi_1-\varphi_2=-\pi/6}$。
\crosscheck 振幅必须取非负根；相位差反序改变符号。

\sourcetest{4.2.5}{95}{103}
\problemtext 相干波源 $S_1,S_2$ 相距 $3\lambda/4$，两源等振幅；在 $S_1$ 左侧各点合波强度为单源4倍，求两源初相差。
\workings 该侧任意点的传播路程差为 $r_2-r_1=3\lambda/4$；强度达到4倍表明同相抵达：$\varphi_2-\varphi_1-2\pi(3/4)=2k\pi$，故 $\boxed{\varphi_2-\varphi_1=3\pi/2\equiv-\pi/2\pmod{2\pi}}$。
\crosscheck “某侧各点”成立的原因是两源在同一直线上，路程差恒定。

\sourcetest{4.2.6}{95}{103}
\problemtext 题4.2.3图给出同一右行波在 $t=0$ 与 $T/4$ 的波形；画出标记点 $P$ 的振动图线。
\workings 从图上 $P$ 在 $t=0$ 经过平衡位置，且波形局部斜率为负；右行波满足 $y_t=-u y_x$，故此时 $v_y>0$。所以可写 $\boxed{y_P(t)=A\sin(2\pi t/T)}$。振动图线从原点以正斜率上升，$T/4$ 达到 $+A$。
\crosscheck 由两幅图的峰值右移 $\lambda/4$ 确认传播方向；与原图 P 的下降空间斜率对应上升时间斜率。

\sourcetest{4.2.7}{95}{103}
\problemtext 题4.2.4图示右行波在 $t=T/4$ 的空间曲线，波速330m/s，振幅0.10m；横轴两相邻波谷距离4m。求波动方程。
\workings $\lambda=4$m，$k=\pi/2$rad/m，$\omega=2\pi(330/4)=165\pi$rad/s。图中该时刻 $x=0$ 过零且波形斜率负，故可选 $\boxed{y=0.10\cos(165\pi t-\pi x/2+\pi)}$（SI）。
\crosscheck 代入 $t=T/4=\pi/(2\omega)$，得 $y(x,T/4)=-0.10\sin(\pi x/2)$，波谷在 $x=1$m，符合原图。

\sourcetest{4.2.8}{95}{103}
\problemtext 右行波 $y=A\cos[2\pi(t/T-x/\lambda)+\varphi]$，分别求 $x=-\lambda$ 质点振动式及以旧坐标 $x=\lambda$ 为原点、反向取轴的新波动式。
\workings 代入 $x=-\lambda$，$\boxed{y_{-\lambda}(t)=A\cos(2\pi t/T+\varphi)}$。若新坐标 $x'=\lambda-x$，则 $x=\lambda-x'$，所以 $\boxed{y=A\cos[2\pi(t/T+x'/\lambda)+\varphi]}$。
\crosscheck 空间坐标轴反向后，相位中空间项的号随之反转；相移一个波长不改变相位模 $2\pi$。

\sourcetest{4.2.9}{95}{103}
\problemtext 左行入射波 $y_i=A\cos[2\pi(t/T+x/\lambda)]$ 在固定端 $x=0$ 发生等幅反射；求反射波式及 $x=2\lambda/3$ 处合振幅。
\workings 固定端位移相消，$\boxed{y_r=-A\cos[2\pi(t/T-x/\lambda)]}$。合波 $y=-2A\sin(2\pi t/T)\sin(2\pi x/\lambda)$，于 $x=2\lambda/3$ 得振幅 $\boxed{A_P=\sqrt3 A}$。
\crosscheck 固定端反射发生 $\pi$ 相变；反射波不是与入射波同向传播。

\sourcetest{4.2.10}{95}{103}
\problemtext 汽车以20m/s速度驶近前方观察者，鸣笛频率1000Hz，空气声速330m/s；求车前方声波波长。
\workings 移动声源朝前发出的相邻波峰间距 $\lambda=(c-v_s)/f=(330-20)/1000$，故 $\boxed{\lambda=0.310\ {\rm m}}$。
\crosscheck 这里车前方受压缩的波长小于静止声源波长0.330m；原题给的1000Hz是鸣笛本征频率。
